\chapter{Uso de la tarjeta con Quartus Prime Pro Edition}

\section{Pasos para cargar proyectos desde \quartus a la tarjeta}

    Para cargar un programa a la tarjeta de desarrollo desde \quartus debemos realizar el siguiente procedimiento con el fin de que se realice correctamente la comunicación mediante protocolo serial desde la MAX 10 al Agilex 7.

    \begin{enumerate}
        \item Conectar el cable USB2.0 tipo B de la tarjeta a la estaci\'on de trabajo.
        \item Revisar que la configuraci\'on de la tarjeta sea la configuraci\'on predeterminada.
        \item Revisa el estatus de los m\'odulos externos: cable MXPM/QSFP/QSFPDD/FMC/DIMM
        \item Encender la tarjeta.
        \item En la estaci\'on de trabajo abrir Quartus Prime Pro Edition
        \item Abrir el proyecto que se va a cargar a la tarjeta de desarrollo.
        \item Establecer las configuraciones necesarias para el dispositivo.
        \item Configurar el programador para que trabaje con el hardware correspondiente a la tarjeta de desarrollo.
        \item Establecer la configuraci\'on del programador e incluir los archivos necesarios.
        \item Ejecutar el proyecto en la tarjeta de desarrollo.
    \end{enumerate}

\subsection{Ejemplo de programa para hacer parpadear un LED}

    Se muestra el siguiente ejemplo que realiza la conmutaci\'on de uno de los LED's de usuario, se muestran los pasos a seguir con el fin de que la comunicaci\'on al chip Agilex 7 se establezca correctamente. 
    
    El proyecto con el cual se realizar\'a este ejemplo se encuentra en el archivo \curl{home/liese2/Proyects/blinky.qpf}, este ser\'a el archivo que nos ayudar\'a a cargar el programa funcional a la tarjeta. Al abrirlo se mostrar\'a la vista inicial del proyecto y en el \texttt{\textcolor{DarkOrchid1}{Project Navigator}} se muestran los archivos que contiene (\cref{fig:screen1}).

    % \begin{figure}[H]
    %     \centering
    %     \includegraphics[width=0.5\linewidth, clip, trim= 0 0 3cm 0]{Images//Parte3/recent.png}
    %     \caption{Proyectos recientes utilizados en \quartus.}
    %     \label{fig:recent}
    % \end{figure}
    \begin{figure}
        \centering
        \includegraphics[width=0.9\linewidth]{Images//Parte3/screen1.png}
        \caption{Vista inicial al abrir el proyecot}
        \label{fig:screen1}
    \end{figure}

    El siguiente paso ser\'a realizar ciertas configuraciones: para ello debemos ir a \texttt{\textcolor{DarkOrchid1}{{Assignments $\rightarrow$ Device $\rightarrow$ Device and Pin Options}}} (\cref{fig:configgen}) para realizar dos configuraciones.
    
    La primera de ellas es en \texttt{\textcolor{DarkOrchid1}{Configuration $\rightarrow$ Configuration Pin Options}}, dentro de este menú se deben marcar las opciones \texttt{\textcolor{DarkOrchid1}{USE PWRMGT\_SCL output}}, \texttt{\textcolor{DarkOrchid1}{USE PWRMGT\_SDA output}} y \texttt{\textcolor{DarkOrchid1}{USE CONF\_DONE output}}. Además de seleccionar \texttt{\textcolor{DarkOrchid1}{SDM\_IO0}}, \texttt{\textcolor{DarkOrchid1}{SDM\_IO12}} y \texttt{\textcolor{DarkOrchid1}{SDM\_IO16}} en las respectivas opciones (\cref{fig:configpin}). 

    \begin{figure}[H]
        \centering
        \begin{subfigure}[c]{0.45\linewidth}
            \centering
            \includegraphics[width=\linewidth]{Images//Parte3/configgen.png}
            \caption{Ventana para la configuraci\'on de la tarjeta y los pines de salida.}
            \label{fig:configgen}
        \end{subfigure}\hspace{.5cm}
        \begin{subfigure}[c]{.45\linewidth}
            \centering
            \includegraphics[width=\linewidth]{Images//Parte3/config1.png}
            \caption{Configuraci\'on para los pines de salida.}
            \label{fig:configpin}
        \end{subfigure}
        \caption{Configuraciones de salida para la tarjeta de desarrollo.}
    \end{figure}


    La siguiente configuraci\'on necesaria es en el apartado de \texttt{\textcolor{DarkOrchid1}{Power Management \& VID}}, en donde se deben establecer los valores de \texttt{\textcolor{DarkOrchid1}{Bus speed mode $\rightarrow$ 100 KHz}}, \texttt{\textcolor{DarkOrchid1}{Slave device type $\rightarrow$ LTC3888}} y \texttt{\textcolor{DarkOrchid1}{PMBus device 0 slave address $\rightarrow$ 62}} (\cref{fig:config2}).

    \textit{Nota: Una vez realizada esta configuraci\'on no es necesario realizarla de nuevo.}

    Al compilar el proyecto tarda unos cinco minutos aproximadamente en completarse y cuando termina lanza una ventana que nos indica que se compil\'o correctamente (\cref{fig:compilado}).

    \begin{figure}[H]
        \centering
        \begin{subfigure}[c]{0.45\linewidth}
            \centering            \includegraphics[width=\linewidth]{Images//Parte3/config2.png}
            \caption{Configuraci\'on para el manejo de energ\'ia.}
            \label{fig:config2}
        \end{subfigure}\hspace{.5cm}
        \begin{subfigure}[c]{.45\linewidth}
            \centering
            \includegraphics[width=\linewidth]{Images//Parte3/compilado.png}
            \caption{Ventana que se lanza cuando se compila el proyecto.}
            \label{fig:compilado}
        \end{subfigure}
        \caption{Configuraci\'on de energ\'ia y compilado.}
    \end{figure}


        Despu\'es de esto nos dirigimos a \texttt{\textcolor{DarkOrchid1}{Tools $\rightarrow$ Programmer}} (\cref{fig:programmer}) y dentro de \texttt{\textcolor{DarkOrchid1}{Hardware Setup}} tenemos que seleccionar el Blaster de la tarjeta, el cual aparece como \texttt{\textcolor{DarkOrchid1}{Agilex I-Series SOC Dev Kit [1-6.2]}} (\cref{fig:blaster}), luego de seleccionarlo simplemente cerramos la ventana haciendo click en \texttt{\textcolor{DarkOrchid1}{Close}} y regresamos al \texttt{\textcolor{DarkOrchid1}{Programmer}}. 

        Tras seleccionar el blaster, en la columna izquierda se habilita la opci\'on \texttt{\textcolor{DarkOrchid1}{Auto Detect}}, en la cual hay que dar clik. Puede que aparezca una advertencia de Quartus Prime (\cref{fig:programmerwarning}), si es as\'i simplemente hay que aceptarla y regresar al programador. Tras realizar la detecci\'on autom\'atica de hardware, en el programador pueden aparecer 2 o 3 chips y el hardware seleccionado de la tarjeta (\cref{fig:programmerdetect}).
        
        es lo siguiente:\begin{figure}[H]
            \centering
            \includegraphics[width=0.8\linewidth]{Images//Parte3/programmer.png}
            \caption{Vista inicial del programador.}
            \label{fig:programmer}
        \end{figure}

        \begin{figure}[H]
        \centering
        \begin{subfigure}[c]{0.45\linewidth}
            \centering            \includegraphics[width=\linewidth]{Images//Parte3/blaster.png}
            \caption{Selecci\'on del blaster.}
            \label{fig:blaster}
        \end{subfigure}\hspace{.5cm}
        \begin{subfigure}[c]{.45\linewidth}
            \centering
            \includegraphics[width=\linewidth]{Images//Parte3/programmerwarning.png}
            \caption{Advertencia que se lanza al ejecutar \texttt{\textcolor{DarkOrchid1}{Auto Detect}}.}
            \label{fig:programmerwarning}
        es lo siguiente:\end{subfigure}
        \caption{Selecci\'on de blaster y error que se puede lanzar al realizar la detecci\'on de hardware.}
    \end{figure}

        \begin{figure}[H]
            \centering
            \includegraphics[width=0.9\linewidth]{Images//Parte3/programmerdetect.png}
            \caption{Cambios en el programador al aplicar \texttt{\textcolor{DarkOrchid1}{Auto Detect}}.}
            \label{fig:programmerdetect}
        \end{figure}
        
        
        Se debe seleccionar el chip con el nombre \texttt{\textcolor{DarkOrchid1}{AGFB027R24C}} y en la columna izquiera dar click en \texttt{\textcolor{DarkOrchid1}{Change File}}, aqu\'i tenems que incluir el archivo \texttt{\textcolor{blue}{blinky.sof}} que se encuentra en el directorio \curl{/home/liese2/Projects/blinky/output_files/} (\cref{fig:programmerfile}) y marcar el checkbox en la columna \texttt{\textcolor{DarkOrchid1}{Program/Configure}} del mismo archivo (\cref{fig:programmerloaded}) y finalmente podemos dar click en \texttt{\textcolor{DarkOrchid1}{Start}} para iniciar el programa en la tarjeta de desarrollo (\cref{fig:programmerstart}).

        \begin{figure}[H]
            \centering
            \includegraphics[width=0.5\linewidth]{Images//Parte3/programmerfile.png}
            \caption{Archivo de salida que se tiene que cargar.}
            \label{fig:programmerfile}
        \end{figure}

        \begin{figure}[H]
            \centering
            \includegraphics[width=0.8\linewidth]{Images//Parte3/programmerloaded.png}
            \caption{Programador con el archivo cargado y la opci\'on de configuraci\'on seleccionada.}
            \label{fig:programmerloaded}
        \end{figure}

        \begin{figure}[H]
            \centering
            \includegraphics[width=0.8\linewidth]{Images//Parte3/programmerstart.png}
            \caption{Ejecuci\'on del programa exitosa.}
            \label{fig:programmerstart}
        \end{figure}

        Cuando se ejecuta el programa en la tarjeta se observa c\'omo el LED \texttt{\textcolor{red}{D0}} comienza a parpadear, lo cual indica que el procedimiento se realiz\'o correctamente. 